\documentclass[11pt, letterpaper]{article}

% --- Idioma y codificación ---
\usepackage[spanish]{babel}
\usepackage[utf8]{inputenc}
\usepackage[T1]{fontenc}
\usepackage{lmodern}

% --- Matemáticas ---
\usepackage{amsmath, amssymb, amsthm}

% --- Formato ---
\usepackage[margin=2.5cm]{geometry}
\usepackage{parskip}
\usepackage{booktabs}
\usepackage{hyperref}
\hypersetup{colorlinks=true, linkcolor=blue!70!black, urlcolor=blue!70!black}

% --- Entornos de teoremas ---
\theoremstyle{definition}
\newtheorem{definition}{Definición}[section]
\theoremstyle{plain}
\newtheorem{theorem}{Teorema}[section]
\newtheorem{lemma}[theorem]{Lema}
\newtheorem{corollary}[theorem]{Corolario}
\theoremstyle{remark}
\newtheorem{remark}{Observación}[section]

\title{%
  T04 --- Demostración formal:\\[4pt]
  Build-heap bottom-up es $O(n)$
}
\author{Curso Linux--C--Rust --- B15 Estructuras de datos en C}
\date{}

\begin{document}
\maketitle
\tableofcontents

% ===================================================================
\section{Modelo}
% ===================================================================

Consideramos un heap binario almacenado en un array de $n$ elementos.
La altura del heap es $h = \lfloor \log_2 n \rfloor$.

\begin{definition}[Convención de niveles]
Numeramos los niveles desde abajo: el nivel~0 son las hojas, el
nivel~$k$ tiene nodos a distancia $k$ del nivel de hojas, y el
nivel~$h$ es la raíz.
\end{definition}

\begin{definition}[Costo de sift-down]
Un sift-down en un nodo de nivel $k$ (medido desde abajo) puede
descender a lo sumo $k$ niveles, con $O(1)$ de trabajo por nivel.
Su costo es $O(k)$.
\end{definition}

\begin{definition}[Algoritmo de Floyd]
Ejecuta sift-down para cada nodo interno, de derecha a izquierda y
de abajo hacia arriba (desde el índice
$\lfloor n/2 \rfloor - 1$ hasta~0).
\end{definition}

% ===================================================================
\section{Demostración 1: conteo de nodos por nivel}
% ===================================================================

\begin{lemma}\label{lem:nodes-per-level}
En un heap completo de $n$ nodos, el número de nodos en el nivel~$k$
(desde abajo) es a lo sumo $\lfloor n / 2^{k+1} \rfloor$.
\end{lemma}

\begin{proof}
En un árbol binario completo con $n$ nodos, el nivel~$d$ (desde arriba)
tiene a lo sumo $2^d$ nodos. El nivel~$d$ desde arriba corresponde al
nivel $k = h - d$ desde abajo, donde $h = \lfloor \log_2 n \rfloor$.

Los nodos en el nivel~$k$ desde abajo tienen profundidad $h - k$
desde arriba, así que hay a lo sumo $2^{h-k}$ nodos.
Como $2^h \leq n < 2^{h+1}$:
\[
  2^{h-k} = \frac{2^h}{2^k} \leq \frac{n}{2^k}
\]

Además, cada nodo de nivel~$k$ es raíz de un subárbol de altura~$k$.
Cada subárbol tiene al menos $2^k$ nodos (si es completo), así que
los subárboles son disjuntos y abarcan como máximo $n$ nodos:
\[
  \text{nodos en nivel } k \leq \left\lfloor \frac{n}{2^k} \right\rfloor
\]

Para el algoritmo de Floyd, solo nos interesan los nodos internos
(nivel $k \geq 1$). El número de nodos internos en nivel~$k$ es a lo
sumo $\lfloor n / 2^{k+1} \rfloor$ (ya que las hojas, nivel~0, ocupan
$\lceil n/2 \rceil$ posiciones).
\end{proof}

% ===================================================================
\section{Demostración 2: la serie
  \texorpdfstring{$\sum k x^k = x/(1-x)^2$}{sum kx\^k = x/(1-x)\^2}}
% ===================================================================

\begin{theorem}\label{thm:series}
Para $|x| < 1$:
\[
  \sum_{k=0}^{\infty} k \cdot x^k = \frac{x}{(1 - x)^2}
\]
\end{theorem}

\begin{proof}
Partimos de la serie geométrica:
\[
  \sum_{k=0}^{\infty} x^k = \frac{1}{1-x} \qquad \text{para } |x| < 1
\]

Derivamos ambos lados respecto a $x$. El lado izquierdo (la derivación
término a término es válida dentro del radio de convergencia):
\[
  \frac{d}{dx} \sum_{k=0}^{\infty} x^k
  = \sum_{k=1}^{\infty} k \cdot x^{k-1}
\]

El lado derecho:
\[
  \frac{d}{dx} \frac{1}{1-x} = \frac{1}{(1-x)^2}
\]

Así:
\[
  \sum_{k=1}^{\infty} k \cdot x^{k-1} = \frac{1}{(1-x)^2}
\]

Multiplicamos ambos lados por $x$:
\[
  \sum_{k=1}^{\infty} k \cdot x^{k} = \frac{x}{(1-x)^2}
\]

Como el término $k=0$ aporta $0 \cdot x^0 = 0$:
\[
  \sum_{k=0}^{\infty} k \cdot x^{k} = \frac{x}{(1-x)^2} \qedhere
\]
\end{proof}

\begin{corollary}\label{cor:series-half}
Evaluando en $x = 1/2$:
\[
  \sum_{k=0}^{\infty} \frac{k}{2^k}
  = \frac{1/2}{(1 - 1/2)^2}
  = \frac{1/2}{1/4} = 2
\]
\end{corollary}

\begin{table}[h]
\centering
\caption{Verificación parcial de la serie $\sum k/2^k$.}
\begin{tabular}{ccc}
\toprule
$k$ & $k/2^k$ & Suma parcial \\
\midrule
0 & 0       & 0      \\
1 & 0.5     & 0.5    \\
2 & 0.5     & 1.0    \\
3 & 0.375   & 1.375  \\
4 & 0.25    & 1.625  \\
5 & 0.15625 & 1.78125 \\
6 & 0.09375 & 1.875  \\
10 & 0.00977 & 1.98047 \\
$\infty$ & --- & \textbf{2} \\
\bottomrule
\end{tabular}
\end{table}

% ===================================================================
\section{Demostración 3: build-heap es
  \texorpdfstring{$O(n)$}{O(n)}}
% ===================================================================

\begin{theorem}\label{thm:build-heap}
El algoritmo build-heap bottom-up (Floyd) ejecuta a lo sumo $n$
comparaciones de sift-down. Su complejidad es $O(n)$.
\end{theorem}

\begin{proof}
El trabajo total $W$ es la suma de los costos de sift-down sobre
todos los nodos internos:
\[
  W = \sum_{k=1}^{h}
      (\text{nodos en nivel } k) \times
      (\text{costo de sift-down en nivel } k)
\]

El costo de sift-down para un nodo en nivel~$k$ es a lo sumo $k$.
El número de nodos en nivel~$k$ es a lo sumo
$\lfloor n / 2^{k+1} \rfloor$ (por el Lema~\ref{lem:nodes-per-level}):
\[
  W \leq \sum_{k=1}^{h}
         \left\lfloor \frac{n}{2^{k+1}} \right\rfloor \cdot k
\]

Acotamos eliminando los pisos ($\lfloor x \rfloor \leq x$):
\[
  W \leq \sum_{k=1}^{h} \frac{n}{2^{k+1}} \cdot k
       = \frac{n}{2} \sum_{k=1}^{h} \frac{k}{2^k}
\]

Acotamos la suma parcial por la serie infinita
(Corolario~\ref{cor:series-half}):
\[
  \sum_{k=1}^{h} \frac{k}{2^k}
  \leq \sum_{k=0}^{\infty} \frac{k}{2^k} = 2
\]

Por lo tanto:
\[
  W \leq \frac{n}{2} \cdot 2 = n \qedhere
\]
\end{proof}

% ===================================================================
\section{Demostración 4: build-heap es
  \texorpdfstring{$\Omega(n)$}{Omega(n)}}
% ===================================================================

\begin{theorem}\label{thm:lower-bound}
Cualquier algoritmo que construya un heap a partir de un array
arbitrario requiere $\Omega(n)$ operaciones.
\end{theorem}

\begin{proof}
Cada uno de los $n$ elementos debe ser examinado al menos una vez:
un elemento no examinado podría ser el mínimo (o máximo) y estar en
una posición incorrecta. Cualquier algoritmo que lo ignore podría
producir un heap inválido.

Formalmente: para $n$ entradas arbitrarias, cualquier algoritmo basado
en comparaciones necesita leer al menos $n$ valores. Esto da una cota
inferior de $\Omega(n)$.

Combinando con el Teorema~\ref{thm:build-heap}:
build-heap bottom-up es $\Theta(n)$.
\end{proof}

% ===================================================================
\section{Demostración 5: build-heap top-down es
  \texorpdfstring{$\Theta(n \log n)$}{Theta(n log n)}}
% ===================================================================

\begin{theorem}\label{thm:top-down}
La construcción top-down (insertar uno a uno con sift-up) tiene
complejidad $\Theta(n \log n)$.
\end{theorem}

\begin{proof}
\emph{Cota superior.} Cada inserción $i$ ejecuta un sift-up de costo
a lo sumo $\lfloor \log_2 i \rfloor$:
\[
  W \leq \sum_{i=1}^{n} \lfloor \log_2 i \rfloor
    \leq \sum_{i=1}^{n} \log_2 n
    = n \log_2 n = O(n \log n)
\]

\emph{Cota inferior.} Los nodos en la segunda mitad del array
($i > n/2$) están a profundidad
$\geq \lfloor \log_2(n/2) \rfloor = \lfloor \log_2 n \rfloor - 1$.
Para una entrada adversarial (por ejemplo, secuencia decreciente en un
max-heap), cada uno de estos $n/2$ nodos ejecuta un sift-up de longitud
$\geq \lfloor \log_2 n \rfloor - 1$:
\[
  W \geq \frac{n}{2} \cdot
    (\lfloor \log_2 n \rfloor - 1) = \Omega(n \log n)
\]

Combinando: $W = \Theta(n \log n)$.
\end{proof}

% ===================================================================
\section{Por qué la diferencia entre top-down y bottom-up}
% ===================================================================

\begin{remark}
La asimetría se resume así:
\begin{itemize}
  \item \textbf{Bottom-up (sift-down):} Muchos nodos ($n/2$ hojas)
    hacen 0~trabajo; pocos nodos (la raíz) hacen $O(\log n)$ trabajo.
    Los pocos nodos caros no compensan los muchos baratos.
  \item \textbf{Top-down (sift-up):} Muchos nodos ($n/2$ en el fondo)
    hacen $O(\log n)$ trabajo; pocos nodos (la raíz) hacen 0~trabajo.
    Los muchos nodos caros dominan.
\end{itemize}
\end{remark}

Formalmente, la diferencia se reduce a qué serie se suma:

\begin{table}[h]
\centering
\begin{tabular}{lcc}
\toprule
\textbf{Método} & \textbf{Serie} & \textbf{Valor} \\
\midrule
Bottom-up & $\displaystyle\sum_{k=0}^{h} k / 2^k$ & $2$ (converge) \\[6pt]
Top-down  & $\displaystyle\sum_{k=0}^{h} k \cdot 2^k / n$
          & $\Theta(\log n)$ por nodo \\
\bottomrule
\end{tabular}
\end{table}

En bottom-up, los pesos ($1/2^k$) decrecen exponencialmente con $k$,
haciendo que la serie converja. En top-down, los pesos ($2^k$) crecen
exponencialmente, haciendo que los términos grandes dominen.

% ===================================================================
\section{Resumen de resultados}
% ===================================================================

\begin{table}[h]
\centering
\begin{tabular}{llp{5.5cm}}
\toprule
\textbf{Resultado} & \textbf{Técnica} & \textbf{Clave} \\
\midrule
Nodos en nivel~$k$: $\leq n/2^{k+1}$
  & Conteo en árbol binario
  & Subárboles disjuntos de altura~$k$ \\
$\sum k/2^k = 2$
  & Derivar serie geométrica
  & $d/dx\,[1/(1{-}x)] = 1/(1{-}x)^2$ \\
Build-heap bottom-up: $O(n)$
  & Sumar trabajo por nivel
  & $W \leq (n/2) \cdot 2 = n$ \\
Build-heap bottom-up: $\Omega(n)$
  & Cota inferior trivial
  & Cada elemento debe examinarse \\
Build-heap top-down: $\Theta(n \log n)$
  & Cotas superior e inferior
  & $n/2$ nodos hacen $\Omega(\log n)$ trabajo \\
\bottomrule
\end{tabular}
\end{table}

\end{document}
